\chapter{Where Horimetrik Lives}
\label{ch:einordnung}

% Register: S44, V13. Sources: Appendix J, experiments_einordnung.py.
% Closing chapter: first the frontiers (Ch. 8), then the map.

A book about a new theory owes its readers an address at the end:
where on the map of mathematics does all of this live? The answer
only assembled itself completely at the very end of the work, and
it is better than we dared to hope. There are, after all, three
possible outcomes for a project like this one. An \emph{unknown}
ground space would be suspicious -- most likely one had simply
searched badly. A \emph{known and complete theory} would be the
end -- one would have repainted old mathematics. The best case is
the third: a known ground space carrying new operations and new
theorems. That is exactly where we have landed.

\section{The inhabitants are old acquaintances}

Prepend a leading zero to the digits of a Hori number,
$e=(0,a_1,\ldots,a_n)$, and every position satisfies $0\le e_j<j$
-- and that is, word for word, the definition of an
\emph{inversion sequence} of length $n{+}1$. Inversion sequences
are in classical bijection with permutations: $e_j$ counts how
many earlier entries of a permutation are larger than the $j$-th.
Hence:

\begin{satz}[Register S44]
$H_n$ is in bijection with the inversion sequences of length
$n{+}1$ and thus with the permutations of $n{+}1$ elements.
Through this lens the structure horizon becomes the permutation
statistic
\[
\sigma(P_a)=n-\min_{e_j>0}\bigl((j{-}1)-e_j\bigr),
\]
and its distribution is known exactly: on each level $1\le r\le n$
there are precisely $r\cdot r!$ elements (the formula holds for
$a\ne0$; zero has $\sigma=0$).
\end{satz}

So the factorial size of the horizons was never a coincidence: the
inhabitants of the Hori worlds \emph{are} permutation codes. This
takes nothing away from the theory -- it says precisely what the
theory is: Horimetrik studies how these well-known codes change
their \emph{value}, their \emph{shape}, and their \emph{identity}
when their space grows. Relocations, projectors, side fidelity,
capacity -- all of that is new on the old ground space.

The statistic itself has a pretty reading: $(j{-}1)-e_j$ is the
distance of a digit from its upper bound. The literature knows the
extreme cases as \emph{tight entries} (digits sitting right at the
top); the structure horizon refines this into a measure: how close
does \emph{any} digit come to its stop? Whether this statistic
corresponds, under the bijection, to a classical permutation
statistic is open (Register V13) -- its distribution $r\cdot r!$
is, at any rate, almost suspiciously clean.

\section{The order: a lattice that already exists}

Comparing digit sequences position by position turns every horizon
into a product of finite chains -- a \emph{distributive lattice}.
In this order the maximally filled shape $M_r$ is an ordinary
element, and the capacity class
$\mathcal F_r=\{P:\sigma(P)\le r\}$ is simply the
\emph{principal ideal} below it: everything that lies under $M_r$.
This is the cleanest abstract description of the structure horizon
we know: $\sigma(P)$ is the first step of the chain
$M_0\le M_1\le M_2\le\cdots$ (with $M_0=0$) whose principal ideal
contains $P$ -- and the extremal principle of
Chapter~\ref{ch:strukturpolynome} becomes, in this language,
almost a tautology. The same positionwise order on inversion
sequences was recently studied as the \emph{middle order} on
permutations, between the weak and the Bruhat order. Whether our
projectors $K_V$, $K_S$ are monotone in this lattice, whether the
commutator is an order defect -- these are fully formulated
follow-up questions addressed to an existing field.

\section{The calculus: old operators, new bookkeeping}

The falling factorials are the house basis of the \emph{finite
operator calculus}: there $\Delta$ plays the role of the
derivative, and the umbral algebra has been studying delta
operators and their basic sequences for decades. Structure
polynomials and the horizon derivative are therefore very well
placed. What Horimetrik adds is a question that is not centre
stage there: not \emph{how} an operator changes a polynomial, but
\emph{how much minimal structure capacity} that costs. The horizon
calculus -- $\Gamma_+$, $\Gamma_\times$, $\Gamma_\Delta$,
$\Gamma_E$, $\Gamma_\circ$, all exact and all attained at the same
canonical object -- is a capacity bookkeeping on top of an
established operator framework. Together with the counting-width
reading from Chapter~\ref{ch:strukturpolynome}, this is the
candidate for a self-contained research paper entirely without the
Hori narrative.

\section{The research programme behind it}

Three further neighbourhoods we have sighted but not entered; let
them be recorded as a programme. \emph{First:} replacing the digit
bounds $0\le a_i\le i$ by arbitrary capacity profiles
$0\le a_i<s_i$ yields the $s$-inversion sequences studied in
combinatorics, together with their geometry (lecture hall
polytopes). An \emph{$s$-Horimetrik} would be a second theory --
which profiles admit a polynomial basis, a value--shape duality, a
side fidelity, is completely open. \emph{Second:} rook theory
writes its polynomials in exactly our basis; which shapes are rook
polynomials of boards, and whether $M_r$ possesses a natural
board, would be a combinatorial route to our product formulas.
\emph{Third:} the counting semantics of the shapes -- coloured
ordered injections -- is the raw material of combinatorial
species; a common language for all the interpretations of this
book might be found there.

\begin{oberflaeche}
The book's conclusion, in three lines: the inhabitants of
Horimetrik are known objects -- permutation codes in a known
lattice, described in the house basis of the calculus of finite
differences. The operations on them -- two ways of growing, two
ways of tidying up, a capacity bookkeeping, a counting width --
are new, and so are the theorems about them. A known ground space
with new operations and new theorems: for a small theory that
began with the reversal of a digit ordering in the dark, that is
probably the best address the map has to offer.
\end{oberflaeche}
